A captura autônoma de alvos não cooperativos por Veículos Aéreos Não Tripulados (VANTs) enfrenta desafios críticos devido à lacuna semântica dos algoritmos de visão clássicos e à latência computacional que desestabiliza a aeronave. Para superar essas limitações, este trabalho desenvolve um sistema de Servovisão Baseada em Posição (PBVS) em tempo real para a aproximação e captura de objetos utilizando uma garra inferior (underside). A arquitetura integra modelos de aprendizado profundo da família YOLO para a detecção robusta e sem marcadores do alvo. Para mitigar oclusões e ruídos causados pelos rotores, um Filtro de Kalman Estendido (EKF) atua no plano da imagem como previsor cinemático. As coordenadas bidimensionais são convertidas para o espaço tridimensional métrico via modelo Pinhole, alimentando um controlador de voo Proporcional-Derivativo (PD) em cascata responsável por compensar a latência visual. O sistema foi modelado em ambiente de simulação Software-in-the-Loop (SITL) utilizando ROS e Gazebo. O projeto estabelece como metas o processamento visual superior a 20 FPS, a redução de 20\% no erro de trajetória (RMSE), a manutenção do sobressinal (overshoot) abaixo de 15\% e uma precisão final de aproximação de 0,30 metros, validando uma solução viável e estável para a manipulação aérea autônoma.